% ============================================================
\chapter{Mouvement Brownien}
\label{ch:brownien}
% ============================================================

En 1827, le botaniste Robert Brown observe au microscope des grains de pollen
suspendus dans l'eau. Ils tremblent, s'agitent, changent de direction sans
cesse --- et personne ne comprend pourquoi. Il faudra attendre 1905 et un
employ\'e du bureau des brevets de Berne, un certain Albert Einstein, pour que
le myst\`ere commence \`a se dissiper~: ces grains sont bombard\'es par des
milliards de mol\'ecules d'eau, et leur trajectoire erratique refl\`ete
l'agitation thermique du fluide. Parall\`element, le math\'ematicien Louis
Bachelier avait d\'ej\`a, dans sa th\`ese de 1900, mod\'elis\'e les fluctuations
boursi\`eres par un processus aux propri\'et\'es remarquablement similaires.

Ce qui nous am\`ene \`a la question fondamentale de ce chapitre~: peut-on
construire un objet math\'ematique rigoureux qui capture cette agitation
al\'eatoire perp\'etuelle~? Un processus continu, qui repart de z\'ero, dont
les incr\'ements ne gardent aucune m\'emoire les uns des autres, et dont
chaque pas suit une loi gaussienne~? La r\'eponse est oui, et cet objet
s'appelle le \emph{mouvement brownien}.

% ------------------------------------------------------------
\section{D\'efinition}
% ------------------------------------------------------------

Commen\c{c}ons par \'enoncer pr\'ecis\'ement ce que nous cherchons \`a construire.
Quatre propri\'et\'es, ni plus ni moins, suffisent \`a caract\'eriser le mouvement
brownien. Chacune encode un aspect physique~: le d\'epart de l'origine,
l'absence de m\'emoire, la normalit\'e des fluctuations, et la continuit\'e
du mouvement.

\begin{definition}[Mouvement brownien standard]
\label{def:brownien}
Un processus stochastique $(B_t)_{t \geq 0}$ \`a valeurs r\'eelles est un
\emph{mouvement brownien standard} (ou \emph{processus de Wiener}) s'il v\'erifie~:
\begin{enumerate}
  \item $B_0 = 0$ p.s.,
  \item $(B_t)$ a des \emph{accroissements ind\'ependants}~: pour tous
    $0 \leq t_1 < t_2 < \cdots < t_n$, les v.a.\
    $B_{t_2} - B_{t_1}, B_{t_3} - B_{t_2}, \ldots, B_{t_n} - B_{t_{n-1}}$
    sont ind\'ependantes,
  \item $(B_t)$ a des \emph{accroissements gaussiens stationnaires}~:
    $B_{t+s} - B_s \sim \mathcal{N}(0, t)$ pour tous $s, t \geq 0$,
  \item les trajectoires $t \mapsto B_t(\omega)$ sont p.s.\ \emph{continues}.
\end{enumerate}
\end{definition}

La premi\`ere condition fixe le point de d\'epart. La deuxi\`eme dit que ce que
fait le processus entre les instants $t_1$ et $t_2$ n'a aucune influence sur
ce qu'il fera entre $t_3$ et $t_4$~: c'est la traduction math\'ematique du
chaos mol\'eculaire. La troisi\`eme impose que chaque incr\'ement suive une loi
normale dont la variance est proportionnelle au temps \'ecoul\'e~: plus on
attend, plus la particule s'\'eloigne, mais en $\sqrt{t}$ seulement. La
quatri\`eme, enfin, demande que les trajectoires soient continues~: la
particule ne \og saute\fg{} jamais.

\begin{remarque}
Ces quatre conditions impliquent que le mouvement brownien est un processus
gaussien centr\'e de fonction de covariance $\Cov(B_s, B_t) = \min(s, t)$.
Pour le voir, il suffit d'\'ecrire $B_s = B_s - B_0$ et
$B_t = B_s + (B_t - B_s)$ pour $s \leq t$, puis d'exploiter l'ind\'ependance
de $B_s$ et $B_t - B_s$.
\end{remarque}

La g\'en\'eralisation \`a plusieurs dimensions est naturelle~: on lance $d$
mouvements browniens ind\'ependants en parall\`ele.

\begin{definition}[Mouvement brownien multidimensionnel]
Un mouvement brownien $d$-dimensionnel est $B_t = (B_t^{(1)}, \ldots, B_t^{(d)})$
o\`u $B^{(1)}, \ldots, B^{(d)}$ sont des mouvements browniens standards ind\'ependants.
\end{definition}

% ------------------------------------------------------------
\section{Existence~: construction de L\'evy--Ciesielski}
% ------------------------------------------------------------

Mais voil\`a o\`u les choses deviennent int\'eressantes~: la d\'efinition ne
garantit pas qu'un tel objet \emph{existe}. Nous avons \'ecrit une liste de
souhaits --- encore faut-il montrer que quelqu'un peut les exaucer. C'est
l\`a qu'intervient la construction de L\'evy--Ciesielski, l'une des plus
\'el\'egantes de l'analyse stochastique.

\begin{theoreme}[Existence du mouvement brownien]
\label{thm:existence_brownien}
Il existe un processus stochastique v\'erifiant les quatre propri\'et\'es de la
d\'efinition~\ref{def:brownien}.
\end{theoreme}

L'id\'ee est de construire le mouvement brownien comme une s\'erie al\'eatoire
dans l'espace des fonctions continues, un peu comme on d\'ecompose un signal
en s\'erie de Fourier --- sauf qu'ici, les coefficients sont des variables
al\'eatoires gaussiennes.

\begin{proof}[Construction de L\'evy--Ciesielski (esquisse)]
On utilise la base de Schauder de $C([0,1])$~: une famille de fonctions
\og chapeau\fg{} dyadiques $s_n$ qui forment une base de cet espace. Soit
$(Z_n)_{n \geq 0}$ une suite de v.a.\ i.i.d.\ $\mathcal{N}(0,1)$. On pose~:
\[
  B_t = \sum_{n=0}^{\infty} Z_n \, s_n(t).
\]
Chaque terme de la s\'erie ajoute un \og tremblement\fg{} \`a une \'echelle de
plus en plus fine. La beaut\'e de la construction est que la s\'erie converge
uniform\'ement p.s.\ dans $C([0,1])$, et la limite est un mouvement brownien
sur $[0,1]$. L'extension \`a $\R_+$ se fait par concat\'enation de copies
ind\'ependantes.
\end{proof}

Norbert Wiener avait d\'ej\`a donn\'e une construction en 1923, mais celle de
L\'evy--Ciesielski, plus tardive, a l'avantage de rendre la convergence
uniforme transparente. Aujourd'hui, d'autres constructions existent
(interpolation gaussienne, th\'eor\`eme de Kolmogorov--Chentsov), mais
celle-ci reste la plus \'eclairante p\'edagogiquement.

% ------------------------------------------------------------
\section{Propri\'et\'es fondamentales}
% ------------------------------------------------------------

Maintenant que nous savons que le mouvement brownien existe, explorons ses
sym\'etries et ses structures cach\'ees. Ces propri\'et\'es ne sont pas de simples
curiosit\'es~: chacune d'entre elles sera un outil de calcul puissant dans
la suite du cours.

Commen\c{c}ons par un fait remarquable~: le mouvement brownien est une
martingale. Intuitivement, cela signifie que la meilleure pr\'ediction de
$B_t$ connaissant le pass\'e est simplement sa valeur actuelle --- le brownien
ne \og d\'erive\fg{} dans aucune direction.

\begin{proposition}[Propri\'et\'es de martingale]
Le mouvement brownien v\'erifie~:
\begin{enumerate}
  \item $(B_t)$ est une martingale par rapport \`a sa filtration naturelle.
  \item $(B_t^2 - t)$ est une martingale.
  \item $\exp(\theta B_t - \theta^2 t/2)$ est une martingale pour tout $\theta \in \R$.
\end{enumerate}
\end{proposition}

La deuxi\`eme propri\'et\'e dit que $B_t^2$ cro\^it en moyenne comme $t$~: c'est
la signature de la diffusion. La troisi\`eme --- la martingale exponentielle
--- est l'outil cl\'e du changement de mesure de Girsanov.

Ce qui nous am\`ene aux \emph{sym\'etries} du mouvement brownien, qui r\'ev\`elent
\`a quel point cet objet est \og autosimilaire\fg. On peut le retoquer, le
dilater, le d\'ecaler dans le temps, voire le retourner --- et on obtient
\`a chaque fois un nouveau mouvement brownien.

\begin{proposition}[Propri\'et\'es d'invariance]
\label{prop:invariance}
Les processus suivants sont encore des mouvements browniens standards~:
\begin{enumerate}
  \item \textbf{Sym\'etrie}~: $(-B_t)_{t \geq 0}$.
  \item \textbf{Scaling}~: $(c^{-1/2} B_{ct})_{t \geq 0}$ pour tout $c > 0$.
  \item \textbf{Translation temporelle}~: $(B_{t+s} - B_s)_{t \geq 0}$ pour tout $s \geq 0$.
  \item \textbf{Inversion temporelle}~: $(tB_{1/t})_{t > 0}$ (avec la convention
    $B_{1/t}\big|_{t=0} = 0$).
\end{enumerate}
\end{proposition}

L'invariance par scaling est particuli\`erement frappante~: elle dit que si
vous zoomez sur une trajectoire brownienne (en ajustant l'\'echelle
verticale en $\sqrt{c}$), vous obtenez exactement la m\^eme loi.
C'est une propri\'et\'e \emph{fractale} --- et nous verrons bient\^ot qu'elle
se manifeste dans la r\'egularit\'e des trajectoires.

\begin{proposition}[Propri\'et\'e de Markov forte]
Pour tout temps d'arr\^et $\tau$ fini p.s., le processus
$(B_{\tau + t} - B_{\tau})_{t \geq 0}$ est un mouvement brownien ind\'ependant
de $\mathcal{F}_{\tau}$.
\end{proposition}

La propri\'et\'e de Markov forte est la version \og am\'elior\'ee\fg{} de la propri\'et\'e
de Markov~: elle dit que l'on peut red\'emarrer le brownien non seulement \`a
un instant d\'eterministe, mais aussi \`a un temps d'arr\^et al\'eatoire. C'est
cette propri\'et\'e qui rend possible l'analyse des temps de premier passage.

% ------------------------------------------------------------
\section{R\'egularit\'e des trajectoires}
% ------------------------------------------------------------

Imaginons que vous tracez une trajectoire brownienne sur un papier, avec un
crayon infiniment fin. Quelle serait la nature de cette courbe~? La r\'eponse
est \`a la fois belle et troublante~: la courbe est continue --- vous ne
l\`everiez jamais le crayon --- mais elle est si tourment\'ee qu'elle n'admet
de tangente \emph{nulle part}. C'est un objet continu mais infiniment
rugueux, une courbe qui d\'efie l'intuition g\'eom\'etrique classique.

\begin{theoreme}[H\"older-continuit\'e]
\label{thm:holder}
Les trajectoires du mouvement brownien sont p.s.\ h\"old\'eriennes d'exposant
$\gamma$ pour tout $\gamma < 1/2$, mais pas pour $\gamma = 1/2$.
\end{theoreme}

L'exposant $1/2$ est le seuil critique~: les trajectoires sont
\og presque\fg{} lipschitziennes d'exposant $1/2$, mais pas tout \`a fait. Ce
\og presque\fg{} a des cons\'equences profondes.

\begin{theoreme}[Non-d\'erivabilit\'e]
\label{thm:non_deriv}
Les trajectoires du mouvement brownien sont p.s.\ \emph{nulle part d\'erivables}.
Plus pr\'ecis\'ement, pour presque tout $\omega$, il n'existe aucun $t \geq 0$ tel que
$\lim_{h \to 0} (B_{t+h}(\omega) - B_t(\omega))/h$ existe (m\^eme comme limite
infinie).
\end{theoreme}

Ce r\'esultat est stup\'efiant quand on y r\'efl\'echit~: nous avons une fonction
continue d\'efinie sur $\R_+$ tout entier qui n'est d\'erivable en aucun point.
Avant Weierstrass (1872), les math\'ematiciens pensaient qu'une telle chose
\'etait impossible. Le mouvement brownien nous offre un exemple
\og naturel\fg{} --- produit par la physique elle-m\^eme, pas par un
contre-exemple artificiel.

Mais les surprises ne s'arr\^etent pas l\`a. La question suivante est~: si les
trajectoires browniennes sont si irr\'eguli\`eres, quelle est leur
\og longueur\fg{}~? Et leur \og \'energie\fg{}~?

\begin{theoreme}[Variation quadratique]
\label{thm:var_quad_brown}
Pour toute suite de partitions $\Pi_n = \{0 = t_0^n < t_1^n < \cdots < t_{k_n}^n = t\}$
de $[0, t]$ avec $\abs{\Pi_n} \to 0$~:
\[
  \sum_{i=0}^{k_n - 1} (B_{t_{i+1}^n} - B_{t_i^n})^2 \xrightarrow{L^2} t.
\]
Autrement dit, $\langle B \rangle_t = t$.
\end{theoreme}

\begin{proof}
La preuve est \'etonnamment courte pour un r\'esultat aussi fondamental. Posons
$Q_n = \sum_{i=0}^{k_n-1} (B_{t_{i+1}} - B_{t_i})^2$. Par lin\'earit\'e de
l'esp\'erance et stationnarit\'e des accroissements~:
$\E[Q_n] = \sum_i (t_{i+1} - t_i) = t$. Il reste \`a montrer que la variance
de $Q_n$ tend vers z\'ero~:
\begin{align*}
  \Var(Q_n) &= \sum_i \Var\bigl((B_{t_{i+1}} - B_{t_i})^2\bigr)
  = \sum_i 2(t_{i+1} - t_i)^2 \\
  &\leq 2 \abs{\Pi_n} \sum_i (t_{i+1} - t_i) = 2 \abs{\Pi_n} \cdot t \to 0.
\end{align*}
Donc $Q_n \to t$ dans $L^2$, ce qui signifie que la somme des carr\'es des
incr\'ements converge vers $t$, quelle que soit la partition choisie.
\end{proof}

\begin{theoreme}[Variation totale infinie]
\label{thm:var_totale}
Les trajectoires du mouvement brownien ont p.s.\ une \emph{variation totale infinie}
sur tout intervalle $[0, t]$ avec $t > 0$~:
\[
  \sup_{\Pi} \sum_i \abs{B_{t_{i+1}} - B_{t_i}} = +\infty \quad \text{p.s.}
\]
\end{theoreme}

\begin{attention}[title=\textbf{Variation quadratique finie, variation totale infinie}]
Voici le paradoxe central qui motive toute la th\'eorie de l'int\'egration
stochastique~: le mouvement brownien a une variation quadratique finie
($= t$) mais une variation totale infinie. Concr\`etement, la \og longueur\fg{}
de la trajectoire est infinie (on ne peut pas la mesurer au sens classique),
mais la \og somme des carr\'es\fg{} des oscillations converge. C'est
pr\'ecis\'ement cette propri\'et\'e qui emp\^eche de d\'efinir l'int\'egrale
stochastique $\int_0^t f \, dB$ au sens de Stieltjes et n\'ecessite la
construction sp\'eciale d'It\^o, que nous verrons au chapitre suivant.
\end{attention}

% ------------------------------------------------------------
\section{Loi du logarithme it\'er\'e}
% ------------------------------------------------------------

Nous savons d\'esormais que le mouvement brownien oscille sans cesse, sans
jamais d\'eriver dans une direction pr\'ecise. Mais \`a quelle \og vitesse\fg{}
s'\'ecarte-t-il de l'origine~? Sa variance cro\^it comme $t$, donc sa
d\'eviation typique est de l'ordre de $\sqrt{t}$. Mais le mouvement brownien
peut-il aller beaucoup plus loin que $\sqrt{t}$, ne serait-ce
qu'occasionnellement~? La loi du logarithme it\'er\'e --- l'un des th\'eor\`emes
les plus fins de la th\'eorie des probabilit\'es --- r\'epond avec une
pr\'ecision chirurgicale.

\begin{theoreme}[Loi du logarithme it\'er\'e]
\label{thm:LIL}
\[
  \limsup_{t \to \infty} \frac{B_t}{\sqrt{2t \log \log t}} = 1 \quad \text{p.s.}
\]
\[
  \liminf_{t \to \infty} \frac{B_t}{\sqrt{2t \log \log t}} = -1 \quad \text{p.s.}
\]
Il existe un r\'esultat analogue en $t \to 0$~:
\[
  \limsup_{t \to 0^+} \frac{B_t}{\sqrt{2t \log \log(1/t)}} = 1 \quad \text{p.s.}
\]
\end{theoreme}

Ce th\'eor\`eme, d\^u \`a Khintchine (1924), est remarquable. Il dit que
l'enveloppe exacte du mouvement brownien est $\pm\sqrt{2t \log \log t}$.
Le brownien d\'epasse cette enveloppe infiniment souvent, mais revient
toujours en dessous. Et la version en $t \to 0$ r\'ev\`ele que m\^eme au
voisinage de l'origine, les oscillations sont d\'ej\`a aussi violentes~: il
n'y a pas d'\og instant calme\fg.

% ------------------------------------------------------------
\section{Principe de r\'eflexion et loi du maximum}
% ------------------------------------------------------------

Changeons de perspective. Supposons que nous lan\c{c}ons un mouvement brownien
et que nous nous demandons~: quand atteindra-t-il pour la premi\`ere fois le
niveau $a > 0$~? Et quel est le maximum atteint sur l'intervalle $[0, t]$~?
Ces questions, fondamentales en finance (barri\`eres d'options) et en
physique (temps de premier passage), trouvent une r\'eponse \'el\'egante
gr\^ace au \emph{principe de r\'eflexion}.

L'id\'ee g\'eom\'etrique est limpide~: une fois que le brownien a touch\'e le
niveau $a$, par la propri\'et\'e de Markov forte, il red\'emarre comme un
brownien neuf. Mais un brownien neuf est sym\'etrique. Donc on peut
\og refl\'eter\fg{} la trajectoire apr\`es le temps d'atteinte, et cette
r\'eflexion a exactement la m\^eme loi que la trajectoire originale.

\begin{theoreme}[Principe de r\'eflexion]
\label{thm:reflexion}
Soit $\tau_a = \inf\{t \geq 0 : B_t = a\}$ pour $a > 0$. Pour tout $b \leq a$~:
\[
  \PP(\tau_a \leq t,\; B_t \leq b) = \PP(B_t \geq 2a - b).
\]
\end{theoreme}

De ce principe d\'ecoulent imm\'ediatement la loi du maximum et celle du
temps de premier passage.

\begin{corollaire}[Loi du maximum]
\label{cor:loi_max}
Soit $M_t = \max_{0 \leq s \leq t} B_s$. Alors~:
\[
  \PP(M_t \geq a) = 2\PP(B_t \geq a) = 2\bigl(1 - \Phi(a/\sqrt{t})\bigr),
  \quad a \geq 0,
\]
o\`u $\Phi$ est la fonction de r\'epartition de $\mathcal{N}(0,1)$.
De plus, $M_t \overset{\mathcal{L}}{=} \abs{B_t}$.
\end{corollaire}

L'identit\'e $M_t \overset{\mathcal{L}}{=} \abs{B_t}$ est surprenante~: le
maximum du brownien sur $[0, t]$ a la m\^eme loi que la valeur absolue du
brownien en $t$. Ce n'est pas une co\"incidence, mais une cons\'equence
profonde de la sym\'etrie du processus.

\begin{corollaire}[Loi de $\tau_a$]
Pour $a > 0$~:
\[
  \PP(\tau_a \leq t) = 2\bigl(1 - \Phi(a/\sqrt{t})\bigr), \quad
  f_{\tau_a}(t) = \frac{a}{\sqrt{2\pi t^3}} \exp\Bigl(-\frac{a^2}{2t}\Bigr),
  \quad t > 0.
\]
En particulier, $\E[\tau_a] = +\infty$.
\end{corollaire}

Ce dernier r\'esultat m\'erite qu'on s'y arr\^ete~: le temps moyen pour
atteindre \emph{n'importe quel} niveau $a \neq 0$ est \emph{infini}. Le
mouvement brownien finit toujours par atteindre $a$ (c'est la r\'ecurrence),
mais il met en moyenne un temps infini \`a le faire. La queue de la
distribution de $\tau_a$ est si lourde ($\sim t^{-3/2}$) que l'esp\'erance
diverge.

% ------------------------------------------------------------
\section{Mouvement brownien et \'equations aux d\'eriv\'ees partielles}
% ------------------------------------------------------------

L'un des ponts les plus profonds des math\'ematiques modernes relie les
probabilit\'es aux \'equations aux d\'eriv\'ees partielles. Ce pont, c'est la
formule de Feynman--Kac, qui dit en substance~: r\'esoudre l'\'equation de la
chaleur, c'est calculer une esp\'erance sur les trajectoires browniennes.
Richard Feynman et Mark Kac l'ont d\'ecouverte ind\'ependamment \`a la fin des
ann\'ees 1940, et elle est devenue l'un des outils les plus puissants de
la physique math\'ematique.

\begin{theoreme}[Formule de Feynman--Kac (cas simple)]
\label{thm:feynman_kac}
Soit $g : \R \to \R$ born\'ee et continue. La fonction
\[
  u(t, x) = \E_x[g(B_T)] = \E[g(x + B_{T-t})], \quad 0 \leq t \leq T,
\]
est l'unique solution born\'ee de l'\'equation de la chaleur~:
\[
  \frac{\partial u}{\partial t} + \frac{1}{2} \frac{\partial^2 u}{\partial x^2} = 0,
  \quad u(T, x) = g(x).
\]
\end{theoreme}

L'\'equation de la chaleur d\'ecrit la diffusion de la temp\'erature dans une
barre m\'etallique. La formule de Feynman--Kac dit que la temp\'erature au
point $x$ et au temps $t$ est exactement la valeur moyenne de la condition
finale $g$, \'evalu\'ee en un point o\`u le brownien issu de $x$ se
retrouverait au temps $T$. La chaleur se propage comme une marche
al\'eatoire continue.

% ------------------------------------------------------------
\section{Temps locaux}
% ------------------------------------------------------------

Terminons ce chapitre par un concept plus subtil~: le \emph{temps local}.
Imaginons que nous \'equipions le niveau $a$ d'un compteur qui s'incr\'emente
chaque fois que le brownien \og passe du temps\fg{} au voisinage de $a$. Ce
compteur, c'est le temps local $L_t^a$~: il mesure la quantit\'e de temps que
le brownien passe autour du point $a$ jusqu'\`a l'instant $t$.

Plus pr\'ecis\'ement, le temps local est la densit\'e de la mesure d'occupation
du brownien.

\begin{definition}[Temps local]
Le \emph{temps local} du mouvement brownien au niveau $a \in \R$ est le processus
$(L_t^a)_{t \geq 0}$ d\'efini par la \emph{formule de densit\'e d'occupation}~:
\[
  \int_0^t f(B_s) \, ds = \int_{-\infty}^{+\infty} f(a) L_t^a \, da
\]
pour toute fonction bor\'elienne $f \geq 0$. \'Equivalemment~:
\[
  L_t^a = \lim_{\varepsilon \to 0} \frac{1}{2\varepsilon}
  \int_0^t \mathbf{1}_{\{\abs{B_s - a} < \varepsilon\}} \, ds.
\]
\end{definition}

Le temps local est l'ingr\'edient cl\'e de la formule de Tanaka, qui est
elle-m\^eme une g\'en\'eralisation de la formule d'It\^o aux fonctions non
lisses. L\`a o\`u It\^o exige $C^2$, Tanaka s'applique \`a la fonction valeur
absolue --- qui n'est pas d\'erivable en z\'ero --- au prix de l'apparition
du temps local.

\begin{theoreme}[Formule de Tanaka]
\label{thm:tanaka}
Pour tout $a \in \R$~:
\[
  \abs{B_t - a} = \abs{a} + \int_0^t \text{sgn}(B_s - a) \, dB_s + L_t^a,
\]
o\`u $\text{sgn}(x) = \mathbf{1}_{\{x > 0\}} - \mathbf{1}_{\{x \leq 0\}}$.
\end{theoreme}

\begin{remarque}
La convention $\text{sgn}(0) = -1$ adopt\'ee ici est arbitraire~: certains
auteurs utilisent $\text{sgn}(0) = 0$ ou $\text{sgn}(0) = +1$. Le choix
n'affecte pas la formule de Tanaka, car l'ensemble $\{s : B_s = a\}$ est
de mesure de Lebesgue nulle p.s., donc la valeur de $\text{sgn}$ en $0$
ne modifie pas l'int\'egrale stochastique $\int_0^t \text{sgn}(B_s - a)\, dB_s$.
\end{remarque}

% ------------------------------------------------------------
\section{Exercices}
% ------------------------------------------------------------

\begin{exercice}
Montrer que le mouvement brownien est p.s.\ non monotone sur tout intervalle.
\end{exercice}

\begin{exercice}
Calculer $\E[B_s B_t]$ pour $s, t \geq 0$ directement \`a partir de la
d\'efinition.
\end{exercice}

\begin{exercice}
V\'erifier que $(t B_{1/t})_{t > 0}$ est un mouvement brownien (avec la
convention $= 0$ en $t = 0$).
\end{exercice}

\begin{exercice}
Utiliser le principe de r\'eflexion pour calculer $\PP(\tau_1 \leq 4)$.
\end{exercice}

\begin{exercice}
Soit $(B_t)$ un mouvement brownien et $f \in C^2(\R)$ avec $f'' \leq 0$.
Montrer que $(f(B_t))$ est une sur-martingale locale.
\end{exercice}

\begin{exercice}
D\'emontrer que $\E[\tau_a] = +\infty$ pour tout $a \neq 0$ en utilisant la
densit\'e de $\tau_a$.
\end{exercice}

\begin{exercice}
Montrer que le pont brownien $X_t = B_t - tB_1$ ($0 \leq t \leq 1$) est un
processus gaussien et calculer sa fonction de covariance. V\'erifier que
$X_0 = X_1 = 0$ p.s.
\end{exercice}
